\documentclass{article}
\usepackage[margin=1in]{geometry}
\usepackage{amsmath,amssymb,amsthm}
\usepackage{booktabs,longtable}
\usepackage{hyperref}

\newtheorem{definition}{Definition}
\newtheorem{theorem}[definition]{Theorem}
\newtheorem{corollary}[definition]{Corollary}
\newtheorem{warning}[definition]{Warning}
\newtheorem{convention}[definition]{Convention}

\newcommand{\source}[1]{\par\noindent\textbf{Source.} #1}
\newcommand{\lean}[1]{\par\noindent\textbf{Lean target.} \texttt{\detokenize{#1}}}
\newcommand{\status}[1]{\par\noindent\textbf{Status.} \texttt{\detokenize{#1}}}
\newcommand{\depends}[1]{\par\noindent\textbf{Depends on.} #1}
\newcommand{\cA}{\mathcal A}
\newcommand{\cB}{\mathcal B}
\newcommand{\cD}{\mathcal D}
\newcommand{\cH}{\mathcal H}
\newcommand{\cM}{\mathcal M}
\newcommand{\cP}{\mathcal P}
\newcommand{\cR}{\mathcal R}
\newcommand{\Ftwo}{\mathbb F_2}

\begin{document}

\section*{Blueprint 013: Final Stabilizer Type Classification}

\subsection*{Scope}

This note promotes the final remaining blueprint block from
\path{agents/2026-06-22-remaining-blueprint-suite-goal.md}.  It supports
\[
  \texttt{CPS-001},\quad
  \texttt{IRANK-001},\quad
  \texttt{CRVN-001},\quad
  \texttt{MTT-001},\quad
  \texttt{TYPE-001},\quad
  \texttt{TYPE-002}.
\]
It is the final source-backed natural-language prerequisite before later Lean
extraction of the canonical stabilizer projection-state type theorem.  It does
not create Lean files or Aristotle packets.

\subsection*{Source and Dependency Basis}

The stabilizer-side theorem is project-local and comes from project proof note
004, especially Theorem 5.1 and Corollaries 5.2--5.3.  The operator-algebra
steps are not reproved here; they are routed through earlier dedicated
blueprints:
\[
\begin{array}{ll}
\texttt{STATE-002},\texttt{STATE-003},\texttt{STATE-004}
  & \text{finite canonical states, quasi-local state, and restriction},\\
\texttt{CPS-001},\texttt{IRANK-001}
  & \text{named projection-state input and internal rank package},\\
\texttt{PROD-001} & \text{product form after adapted symplectic reduction},\\
\texttt{GNS-001},\texttt{GNS-004},\texttt{GNS-003},\texttt{VN-001},\texttt{VN-003}
  & \text{canonical-state GNS data and represented-region transport},\\
\texttt{CRVN-001}
  & \text{canonical-state represented region von Neumann algebra},\\
\texttt{TRACE-001},\texttt{TYPE-DEF-001},\texttt{TYPE-ALT-001}
  & \text{trace criteria and direct factor type predicates},\\
\texttt{MFH-001},\texttt{ITP-PURE-001},\texttt{ITP-TRACE-001},\texttt{TENSOR-001},
\texttt{MTT-001}
  & \text{pure--defect model factors and tensor type table}.
\end{array}
\]
The projection-defined strengthening needed for the final theorem is named
below as \(\mathrm{CPS}_{\Gamma,S'}\) and carried by \texttt{CPS-001}.

\begin{convention}[Canonical projection-state scope]
Let \(\Gamma\) be a finite or countably infinite qubit region.  Let \(S'\) be a
consistent stabilizer extension with binary stabilizer label space \(W'\), and
let
\[
  W'_\Gamma = W'\cap V_\Gamma
\]
be its internal region stabilizer space.  The state under consideration is the
canonical stabilizer projection state \(\omega_{S'}\), restricted to the region
quasi-local algebra \(\cA_\Gamma\).  This note does not classify arbitrary
logical ground states, arbitrary non-canonical states, or states obtained by
unproved quasi-local/gLU normal forms.
\source{Project proof note 004, Convention 2.3, Lemma 2.4, Proposition 4.1,
Theorem 5.1, and Warning 6.2, local PDF
\path{../runs/2026-06-18-stabilizer-algebra-classification/artifacts/004-internal-stabilizer-reduction-type-classification-R1.pdf},
PDF pp.~2--7; \path{blueprints/009-internal-stabilizer-reduction.tex}.}
\status{Scope convention for \texttt{TYPE-001} and \texttt{TYPE-002}.}
\end{convention}

\subsection*{Named Inputs for the Capstone Theorem}

\begin{definition}[ProjectionDefinedCanonicalStateInput, ID: CPS-001, project proof note 004 Convention 2.3, Lemma 2.4, Proposition 4.1, and Theorem 5.1, PDF pp.~2--6]
The named input
\(\mathrm{CPS}_{\Gamma,S'}\) means the following projection-defined
canonical-state package.
For every finite \(F\subseteq\Gamma\), there is a finite code projection
\(q_F\in\cA_F\) onto the simultaneous eigenspace cut out by the internal
stabilizer labels \(W'\cap V_F\) and their character values.  The finite
state is the normalized projection trace
\[
  \omega_F(a)=\frac{\operatorname{Tr}_F(q_Fa)}{\operatorname{Tr}_F(q_F)},
  \qquad a\in\cA_F,
\]
and this projection-state formula has the Pauli coefficient rule of
\texttt{STATE-002}.  The finite states are compatible under inclusions, extend
uniquely to the quasi-local state \(\omega_{S'}:\cA_\Gamma\to\mathbb C\) of
\texttt{STATE-003}, and the restriction to any region is controlled by the
internal stabilizer space as in \texttt{STATE-004}.
\source{Project proof note 004, Convention 2.3, Lemma 2.4, Proposition 4.1,
and Theorem 5.1, PDF pp.~2--6; \path{blueprints/002-canonical-state.tex};
\path{blueprints/008-stabilizer-local-state-foundations.tex}, targets
\texttt{STATE-002}, \texttt{STATE-003}, and \texttt{STATE-004}.}
\depends{\texttt{STATE-002}, \texttt{STATE-003}, \texttt{STATE-004}.  The
finite code-projection construction is recorded in
\path{blueprints/008-stabilizer-local-state-foundations.tex} as the
projection-state layer supporting \texttt{STATE-002}.}
\lean{Formalization.Stabilizer.RegionType.ProjectionDefinedCanonicalStateInput}
\status{blueprint_note_complete}
\end{definition}

\begin{proof}[Well-formedness]
The finite projection \(q_F\) lives in the finite matrix algebra
\(\cA_F\), so the normalized trace formula is a state once
\(\operatorname{Tr}_F(q_F)>0\).  Compatibility and the quasi-local extension
are precisely the data needed to make \(\omega_{S'}\) independent of the
finite stage.  The definition records this as a named input rather than using
the phrase ``canonical projection state'' as an unexpanded assumption.
\end{proof}

\begin{definition}[InternalRanks, ID: IRANK-001, project proof note 004 Definition 3.3 and Theorem 5.1, PDF pp.~4--6]
Choose an algebraic adapted symplectic basis
\[
  \{z_i,x_i:i\in I\}
\]
for \(V_\Gamma\) such that
\[
  W'_\Gamma=\operatorname{span}_{\Ftwo}\{z_i:i\in P\}
\]
for some \(P\subseteq I\).  Define the internal fixed rank and internal defect
rank by
\[
  p:=|P|=\dim_{\Ftwo}W'_\Gamma,\qquad
  r:=|I\setminus P|.
\]
The ranks are finite or countably infinite.
\source{Project proof note 004, Definition 3.3 and Theorem 5.1, PDF
pp.~4--6; \path{blueprints/009-internal-stabilizer-reduction.tex},
targets \texttt{SYMP-001}, \texttt{SYMP-002}, and \texttt{PROD-001}.}
\depends{\texttt{SYMP-001}, \texttt{SYMP-002}, \texttt{PROD-001}.}
\lean{Formalization.Stabilizer.InternalRanks}
\status{blueprint_note_complete}
\end{definition}

\begin{definition}[CanonicalStateRegionVN, ID: CRVN-001, Takesaki I Theorem I.9.14 and Definition II.3.2, PDF pp.~59--61, 92]
Assume \(\mathrm{CPS}_{\Gamma,S'}\).  Applying \texttt{GNS-001} to the
state \(\omega_{S'}|_{\cA_\Gamma}\) gives the canonical-state GNS triple
\[
  (\cH_{\omega_{S'}},\pi_{\omega_{S'}},\Omega_{\omega_{S'}})
\]
with
\[
  \omega_{S'}(a)=
  \langle \pi_{\omega_{S'}}(a)\Omega_{\omega_{S'}},
          \Omega_{\omega_{S'}}\rangle .
\]
The represented region von Neumann algebra of the canonical projection state
is then
\[
  \cM_\Gamma(S') :=
  \pi_{\omega_{S'}}(\cA_\Gamma)''\subseteq\cB(\cH_{\omega_{S'}}).
\]
\source{Project proof note 004, Theorem 5.1, PDF pp.~5--6;
\path{blueprints/010-gns-represented-region-layer.tex}, targets
\texttt{GNS-001}, \texttt{GNS-003}, \texttt{GNS-004}, \texttt{VN-001}, and
\texttt{VN-003};
Takesaki I, Theorem I.9.14, PDF pp.~59--61, Definition I.9.15, PDF p.~61,
Definition II.3.2, PDF p.~92, and Theorem II.3.9, PDF pp.~94--95.}
\depends{\texttt{CPS-001}, \texttt{GNS-001}, \texttt{GNS-003},
\texttt{GNS-004}, \texttt{VN-001}, and \texttt{VN-003}.}
\lean{Formalization.Stabilizer.RegionType.regionVN}
\status{blueprint_note_complete}
\end{definition}

\begin{proof}[Well-formedness]
The named input \(\mathrm{CPS}_{\Gamma,S'}\) supplies a state on the
quasi-local algebra, so the GNS theorem applies.  The region algebra
\(\cA_\Gamma\) is represented on the resulting Hilbert space, and the double
commutant construction in \texttt{VN-001} turns its represented image into a
concrete von Neumann algebra.  The product-state identification and spatial
transport used later are not part of the definition; they enter through
\texttt{GNS-003} and \texttt{VN-003}.
\end{proof}

\begin{definition}[RegionModelTraceTypeInput, ID: MTT-001]
For \(P,R\subseteq I\), the named input
\(\mathrm{MTT}_{P,R}\) means that the product model
\[
  \cP_P\bar\otimes\cD_R
\]
is interpreted using the direct factor type predicates of
\texttt{TYPE-DEF-001}, the trace and semifiniteness criteria of
\texttt{TRACE-001}, and the four-case model-factor theorem
\texttt{TENSOR-001}.  If the Blackadar-backed model-factor proofs inside
\texttt{TENSOR-001} have not yet been formalized, the only permitted extra
assumptions are the named hypotheses
\(\mathrm{MFH}_{\mathrm{factor}}\),
\(\mathrm{MFH}_{\mathrm{tracial}\,M_2}\), and
\(\mathrm{MFH}_{\cR\otimes\cB}\) from
\path{blueprints/012-infinite-product-type-calculus.tex}.

The GNS vector state on \(\cM_\Gamma(S')\) is normal by construction, but it
is not assumed to be tracial in general.  The finite or semifinite traces used
to recognize the displayed types are the model traces supplied by
\texttt{ITP-TRACE-001}, \texttt{TENSOR-001}, and \texttt{TRACE-001}, then
transported across the spatial isomorphism of \texttt{VN-003}.
\source{\path{blueprints/011-von-neumann-type-theory.tex}, targets
\texttt{TRACE-001}, \texttt{TYPE-DEF-001}, and \texttt{TYPE-ALT-001};
\path{blueprints/012-infinite-product-type-calculus.tex}, targets
\texttt{MFH-001}, \texttt{ITP-PURE-001}, \texttt{ITP-TRACE-001}, and
\texttt{TENSOR-001}.}
\depends{\texttt{TRACE-001}, \texttt{TYPE-DEF-001}, \texttt{TYPE-ALT-001},
\texttt{MFH-001} if deferred, \texttt{ITP-PURE-001},
\texttt{ITP-TRACE-001}, and \texttt{TENSOR-001}.}
\lean{Formalization.Stabilizer.RegionType.ModelTraceTypeInput}
\status{named_input_for_trace_and_type_criteria}
\end{definition}

\begin{proof}[Well-formedness]
The input separates three pieces that are otherwise easy to conflate: normality
of the GNS vector state, existence of model traces, and direct type predicates.
All traces and type predicates are attached to concrete represented von
Neumann algebras, and spatial transport preserves the relevant algebraic
properties.  Thus the final theorem can cite \(\mathrm{MTT}_{P,R}\) without
asserting that the canonical GNS vector state itself is a trace.
\end{proof}

\subsection*{Final Type Classification}

\begin{theorem}[RegionTypeClassification, ID: TYPE-001, project proof note 004 Theorem 5.1, PDF pp.~5--6]
Let \(\Gamma\), \(S'\), \(\omega_{S'}\), \(p\), \(r\), and
\(\cM_\Gamma(S')\) be as above.  Assume the named inputs
\(\mathrm{CPS}_{\Gamma,S'}\) and \(\mathrm{MTT}_{P,I\setminus P}\), and assume
the adapted product-form and GNS transport hypotheses listed in the dependency
line below.  Then \(\cM_\Gamma(S')\) satisfies the direct type predicates of
\texttt{TYPE-DEF-001} as follows:
\[
\begin{array}{c|c|c}
\text{internal fixed rank }p & \text{internal defect rank }r
  & \text{type of }\cM_\Gamma(S')\\
\midrule
p<\infty & r<\infty & \text{finite type } I\\
p=\infty & r<\infty & \text{type } I_\infty\\
p<\infty & r=\infty & \text{hyperfinite type } II_1\\
p=\infty & r=\infty & \text{type } II_\infty .
\end{array}
\]
\source{Project proof note 004, Theorem 5.1, PDF pp.~5--6.  Stabilizer
dependencies: \path{blueprints/008-stabilizer-local-state-foundations.tex} for
region restriction and crossing-stabilizer invisibility;
\path{blueprints/009-internal-stabilizer-reduction.tex} for adapted
symplectic reduction and product form.  Operator-algebra dependencies:
\path{blueprints/010-gns-represented-region-layer.tex} for \texttt{GNS-004}
canonical projection-code GNS specialization and GNS transport;
\path{blueprints/011-von-neumann-type-theory.tex} for type vocabulary; and
\path{blueprints/012-infinite-product-type-calculus.tex} for the precise
pure--defect model-factor type table and named conditional model-factor
hypotheses.}
\depends{\texttt{CPS-001}, \texttt{IRANK-001}, \texttt{CRVN-001},
\texttt{STATE-002},
\texttt{STATE-003}, \texttt{STATE-004}, \texttt{SYMP-001},
\texttt{SYMP-002}, \texttt{PROD-001}, \texttt{GNS-001},
\texttt{GNS-004}, \texttt{GNS-003}, \texttt{VN-001}, \texttt{VN-003},
\texttt{TRACE-001},
\texttt{ITP-PURE-001}, \texttt{ITP-TRACE-001}, \texttt{TENSOR-001},
\texttt{TYPE-DEF-001}, \texttt{TYPE-ALT-001}, \texttt{MATRIX-001}, and
\texttt{MTT-001}.}
\lean{Formalization.Stabilizer.RegionType.type_classification}
\status{blueprint_note_complete}
\end{theorem}

\begin{proof}
By \texttt{CPS-001}, equivalently by
\(\mathrm{CPS}_{\Gamma,S'}\), the state is the quasi-local extension of
finite normalized code-projection states.  By \texttt{STATE-004}, the
coefficients of
\(\omega_{S'}|_{\cA_\Gamma}\) depend only on the internal stabilizer subspace
\(W'_\Gamma\); crossing stabilizers affect the local type only through
internal products that lie in \(W'_\Gamma\).  By \texttt{SYMP-001} and
\texttt{SYMP-002}, an algebraic adapted symplectic basis transports the region
Pauli algebra to a formal product qubit algebra with
\[
  W'_\Gamma=\operatorname{span}_{\Ftwo}\{z_i:i\in P\}.
\]
By \texttt{PROD-001}, the transported state is the product of pure
\(Z\)-eigenstates over \(P\) and normalized traces over \(I\setminus P\):
\[
  \bigotimes_{i\in P}\omega_{Z,\epsilon_i}
  \otimes
  \bigotimes_{i\in I\setminus P}\operatorname{tr}_2 .
\]
The signs \(\epsilon_i\in\{\pm1\}\) record the stabilizer character.  They
change only the one-site pure vector states and are removed by one-site unitary
conjugacies, so they do not change the represented model factor type.

The canonical-state GNS specialization in \texttt{GNS-004} supplies
\((\cH_{\omega_{S'}},\pi_{\omega_{S'}},\Omega_{\omega_{S'}})\), and
\texttt{VN-001} defines
\(\cM_\Gamma(S')=\pi_{\omega_{S'}}(\cA_\Gamma)''\).  By \texttt{GNS-003} and
\texttt{VN-003}, state transport by the adapted \(*\)-isomorphism gives a
spatial isomorphism between \(\cM_\Gamma(S')\) and the represented
pure--defect product factor
\[
  \cP_P\bar\otimes\cD_{I\setminus P},
\]
where \(\cP_P\) is the pure-mode product factor and \(\cD_{I\setminus P}\) is
the tracial-defect product factor of
\path{blueprints/012-infinite-product-type-calculus.tex}.

The named input \texttt{MTT-001}, written
\(\mathrm{MTT}_{P,I\setminus P}\), says that the product model is classified
by \texttt{TENSOR-001}, with any deferred Blackadar model-factor claims
exposed only through \texttt{MFH-001}.  Applying \texttt{TENSOR-001} gives
exactly the four displayed cases.  The traces used
to identify the finite and semifinite cases are the model traces from
\texttt{TRACE-001} and \texttt{TENSOR-001}; the proof does not require the
canonical GNS vector state to be tracial.  Spatial isomorphism preserves the
factor type predicates from \texttt{TYPE-DEF-001}, so the same direct type
table holds for \(\cM_\Gamma(S')\).
\end{proof}

\begin{corollary}[NoTypeIIIRegion, ID: TYPE-002, project proof note 004 Theorem 5.1, PDF pp.~5--6]
No canonical projection state of a Pauli stabilizer extension produces a type
III represented region algebra under this finite/countable internal-stabilizer
classification.
\source{Project proof note 004, Theorem 5.1, PDF pp.~5--6; Takesaki I,
Theorem V.1.19, Corollary V.1.20, and Definition V.1.21, PDF pp.~316--317,
for factor alternatives and semifiniteness; Blackadar III.3.1.4, PDF p.~302;
III.3.1.6, PDF p.~303; and III.3.4.2, PDF p.~317, for the \(II_1\) and \(II_\infty\) model
identifications used in \texttt{TENSOR-001}.}
\depends{\texttt{TYPE-001}, \texttt{TRACE-001}, \texttt{TYPE-DEF-001},
\texttt{TYPE-ALT-001}, \texttt{TENSOR-001}, and \texttt{MTT-001}.}
\lean{Formalization.Stabilizer.RegionType.no_typeIII}
\status{blueprint_note_complete}
\end{corollary}

\begin{proof}
The four cases of \texttt{TYPE-001} are finite type I, type \(I_\infty\),
hyperfinite type \(II_1\), and type \(II_\infty\), stated using the direct
predicates of \texttt{TYPE-DEF-001}.  By the factor type alternatives in
\texttt{TYPE-ALT-001}, a factor satisfying any one of these four predicates is
not type III.  Equivalently, \(\mathrm{MTT}_{P,I\setminus P}\) supplies finite
or faithful semifinite normal model traces in all four cases, and
\texttt{TRACE-001} identifies these as semifinite cases rather than the
non-semifinite type III alternative.
\end{proof}

\begin{warning}[No gLU or arbitrary-state conclusion]
The proof uses an algebraic adapted symplectic basis and the projection-defined
canonical state whose coefficients are computed by \(\mathrm{CPS}_{\Gamma,S'}\).
It does not assert that the basis change is implemented by a geometrically
local Clifford circuit, finite-depth circuit, gLU, or quasi-local automorphism.
It also does not classify arbitrary logical ground states or arbitrary states
extending the same stabilizer constraints.
\source{Project proof note 004, Warning 6.2, PDF p.~7; current goal file
\path{agents/2026-06-22-remaining-blueprint-suite-goal.md}, Non-Goals.}
\status{Scope caveat for later Lean theorem statements.}
\end{warning}

\subsection*{Lean Extraction Notes}

\begin{longtable}{@{}p{0.16\linewidth}p{0.34\linewidth}p{0.40\linewidth}@{}}
\toprule
Target & Intended file & Notes \\
\midrule
\endhead
\texttt{CPS-001} &
\texttt{Formalization/Stabilizer/RegionTypeInputs.lean} &
Named projection-defined canonical-state input package; use before stating the
capstone theorem. \\
\texttt{IRANK-001} &
\texttt{Formalization/Stabilizer/RegionTypeInputs.lean} &
Internal fixed and defect ranks extracted from the adapted symplectic product
form. \\
\texttt{CRVN-001} &
\texttt{Formalization/Stabilizer/RegionTypeInputs.lean} &
Canonical-state GNS represented region algebra using \texttt{GNS-004} and
\texttt{VN-001}. \\
\texttt{MTT-001} &
\texttt{Formalization/Stabilizer/RegionTypeInputs.lean} &
Named model trace/type input package, with \texttt{MFH-001} explicit whenever
the Blackadar model-factor proofs are deferred. \\
\texttt{TYPE-001} &
\texttt{Formalization/Stabilizer/TypeRestrictions.lean} &
State the theorem for canonical projection states, finite/countable regions,
and the internal ranks \(p,r\).  Expose
\(\mathrm{CPS}_{\Gamma,S'}\), \(\mathrm{MTT}_{P,I\setminus P}\), and any
\(\mathrm{MFH}\) model-factor assumptions explicitly until the corresponding
earlier Lean targets exist. \\
\texttt{TYPE-002} &
\texttt{Formalization/Stabilizer/TypeRestrictions.lean} &
Should be a short corollary of \texttt{TYPE-001} plus the project-local factor
type alternatives and trace criteria.  Do not generalize it to arbitrary
logical states. \\
\bottomrule
\end{longtable}

\begin{thebibliography}{9}
\bibitem{TakesakiI}
M. Takesaki, \emph{Theory of Operator Algebras I}.  Local source:
\path{../sources/takesaki_1.pdf}.

\bibitem{Blackadar}
B. Blackadar, \emph{Operator Algebras}.  Local source:
\path{../sources/blackadar.pdf}.

\bibitem{ProjectNote004}
Project proof note 004,
\emph{Internal Stabilizer Reduction and Projection-State Type Classification}.
Local source:
\path{../runs/2026-06-18-stabilizer-algebra-classification/artifacts/004-internal-stabilizer-reduction-type-classification-R1.pdf}.
\end{thebibliography}

\end{document}
